\section{Introduction}

Human gait provides an accessible signal for assessing balance, coordination, strength, and mobility~\cite{perry1992gait,baker2006gait,muro2014gait}. Advances in video-based pose estimation allow ordinary walking videos to be converted into skeleton trajectories without specialized motion-capture systems~\cite{cao2017openpose,lugaresi2019mediapipe,bazarevsky2020blazepose}. Clinical reviews describe how this low-cost route has been used across neurological and rehabilitation settings~\cite{stenum2021applications,stenum2024clinical}. Pose trajectories provide a compact, anatomically indexed representation from which interpretable measurements of motion, symmetry, stability, and variability can be derived~\cite{winter2009biomechanics,schwartz2008gdi}.

Most pose-based clinical gait methods are task-specific supervised systems that map skeleton sequences to gait classes, severity levels, or clinical scores~\cite{khalil2022ataxia,deng2024parkinsonism}. Clinical studies have related quantitative gait measurements to motor severity and function in cerebral palsy, Parkinsonian impairment, and ataxia~\cite{kidzinski2020video,lu2020updrs,khalil2022ataxia}. Related visual studies span heterogeneous cohorts, acquisition settings, and label definitions~\cite{jun2020pathological,deng2024parkinsonism,khalil2022ataxia}, complicating direct cross-resource comparisons. Thus, cross-condition evaluation is not itself the novelty of this work. Instead, we ask whether LLM predictions are grounded in supplied pose-derived evidence. This distinction is important because many gait resources contain numerous clips from relatively few participants, requiring participant-aware interpretation.

Pose evidence can be presented as raw coordinates, deterministic motion summaries, or biomechanical features. Poor performance may therefore reflect limited task-relevant signal, difficulty interpreting the representation, or failure to map informative evidence to dataset-specific labels. Accuracy alone cannot distinguish these cases, particularly for imbalanced ordinal and multi-class tasks~\cite{opitz2019macro,sokolova2009performance}. Evaluation should therefore compare LLMs with label-prior baselines and simple classifiers trained on the same evidence.

We introduce \textbf{GaitEval}, a multi-condition benchmark in which \emph{gait reasoning} means grounding dataset-specified label predictions in pose-derived evidence, rather than open-ended diagnosis or clinical interpretation. The benchmark comprises six datasets: ASD-Kinect, ATAXIA, CARE-PD, GMFCS, PARKINSONISM, and PATHOGAIT. Each retains its population provenance and native label space. GaitEval introduces no new clinical recordings; it provides controlled representations, evidence-use tests, and same-feature diagnostic references (Fig.~\ref{fig:method_pipeline}).

The evaluation follows the original evidence ladder: zero-shot prompting, few-shot examples, and LoRA adaptation, with transfer limited to the tested ATAXIA--PATHOGAIT pair~\cite{brown2020fewshot,hu2022lora}. The common six-dataset analysis uses Qwen2.5-7B-Instruct, subject-disjoint partitions, and capture-quality-free features. Additional model, view, and adaptation measurements retain their original protocol labels and are never substituted for missing common-protocol experiments.

Our main contributions are:
\begin{itemize}[leftmargin=*,nosep]
\item We introduce GaitEval, a six-dataset benchmark for testing whether LLM predictions are grounded in pose-derived gait evidence across heterogeneous resources, representations, and label spaces.
\item We provide a baseline-anchored protocol that separates limited input signal from failed evidence use by comparing controlled LLM inputs with label-prior and linear feature-based references.
\item We evaluate an evidence ladder from zero-shot prompting to few-shot adaptation and LoRA fine-tuning within its measured coverage, showing in-distribution recovery on selected tasks but no reliable transfer on the tested compatible dataset pair.
\end{itemize}
